\section{Introduction and mathematical motivation}
\label{sec:introduction}

The CQTS project investigates whether suitably homotopical topological quantum
transformations can acquire an executable realization as dependent-type
transport in Cubical Agda.  The motivating perspective is the programme of
Myers--Sati--Schreiber \cite{mss}.  Its later stages require precise higher
types, dependent families over them, loops, and a computation analysis of
transport.  None of those later constructions follows merely from naming a
familiar group or classifying space.

This paper records a completed smaller layer needed to formulate one possible
such direction.  Jackson Brough's constructive HoTT real numbers provide a
checked real-number object $\mathbb{R}_{\mathrm H}$.  From it the local
development constructs
\begin{equation}
  \CH\coloneqq\mathbb{R}_{\mathrm H}\times\mathbb{R}_{\mathrm H},
\end{equation}
proves that it is a commutative ring, takes its algebraic units $\CHunits$,
and applies the existing generic theory of left torsors.  The result is the
pointed torsor-moduli type
\begin{equation}
 (\BG\CHunits,b_{\CHunits})\quad\text{with}\quad
 \Loop(\BG\CHunits,b_{\CHunits})\simeq\CHunits .
 \label{eq:main-result}
\end{equation}

The notation $\BG\CHunits$ is a name for this torsor-moduli construction, not
a claim that a separate conventional geometric classifying-space model has
been compared with it.  Equation~\eqref{eq:main-result} is the exact
type-equivalence endpoint of the formalization.  It is neither a newly
packaged group isomorphism nor a construction of a quantum gate.

The argument is layered.  Section~\ref{sec:complex} gives the constructive
complex ring; Section~\ref{sec:units} its algebraic unit group; and
Sections~\ref{sec:torsors}--\ref{sec:loops} specialize the existing
univalent torsor delooping and trace its loop computation.  The provenance
map and reproducibility record are included in the appendices.
